\documentclass[10pt]{article}
\usepackage[letterpaper,margin=0.55in]{geometry}
\usepackage{booktabs}
\usepackage{amsmath}
\usepackage{float}
\begin{document}
\begin{table}[H]
\centering
\caption{Accuracy, peak inference memory, and checkpoint storage for serving four forecasting horizons. Accuracy is the macro average of the 28 Main I dataset--horizon cells. Memory and storage are macro-averaged over seven datasets. Lower is better.}
\label{tab:iscf-bsca-efficiency}
\small
\setlength{\tabcolsep}{7pt}
\begin{tabular}{lrrrr}
\toprule
System & MSE & MAE & Peak memory (MiB) & Checkpoints (MiB) \\
\midrule
\textbf{ISCF-BSCA (ours)} & \textbf{0.261} & \textbf{0.306} & 38.817 & \underline{17.677} \\
TimeAlign & \underline{0.274} & \underline{0.314} & 109.102 & 95.433 \\
QDF & 0.288 & 0.331 & \underline{31.939} & 20.381 \\
AMD & 0.282 & 0.328 & 238.035 & 221.150 \\
SimpleTM & 0.293 & 0.333 & \textbf{14.841} & \textbf{2.835} \\
\bottomrule
\end{tabular}
\vspace{2pt}
\parbox{0.98\linewidth}{\footnotesize ISCF-BSCA stores and serves one unified checkpoint. TimeAlign, QDF, AMD, and SimpleTM sum the four native checkpoints for $H\in\{96,192,336,720\}$. Peak memory is measured on an exclusive RTX 3090 in FP32 with batch size 1 and all service checkpoints resident: one $H=720$ forward plus prefix views for ISCF-BSCA, and four sequential native forwards for each horizon-specific family. SimpleTM resource cost uses repeat 0 as one deployable, non-ensemble instance per horizon, whereas its frozen Main I accuracy follows the table's multi-run summary. Synthetic standardized inputs are used; no test loader or labels are accessed. Best results are bold and second-best results are underlined.}
\end{table}

\end{document}
